% ============================================================
% QHAWARINA — Reporte Trimestral
% Template: reporte_trimestral.tex
% Placeholders: PLACEHOLDER replaced by generate_reports.py
% ============================================================
\documentclass[10pt,a4paper]{article}

\usepackage[margin=2cm,top=2.5cm,bottom=2.5cm]{geometry}
\usepackage{mathpazo}
\usepackage{xcolor}
\usepackage{booktabs}
\usepackage{tabularx}
\usepackage{tcolorbox}
\usepackage{fancyhdr}
\usepackage{titlesec}
\usepackage{graphicx}
\usepackage{hyperref}
\usepackage[spanish]{babel}
\usepackage{ragged2e}
\usepackage{microtype}
\usepackage{array}
\usepackage{colortbl}
\usepackage{longtable}
\usepackage{multicol}
\usepackage{enumitem}

% ------------------------------------------------------------------
% Colour definitions
% ------------------------------------------------------------------
\definecolor{terra}{HTML}{C65D3E}
\definecolor{ink}{HTML}{2D3142}
\definecolor{ink3}{HTML}{8D99AE}
\definecolor{teal}{HTML}{2A9D8F}
\definecolor{amber}{HTML}{E0A458}
\definecolor{red}{HTML}{9B2226}
\definecolor{bg}{HTML}{FAF8F4}
\definecolor{border}{HTML}{E8E4DF}

\pagecolor{bg}
\color{ink}

% ------------------------------------------------------------------
% Hyperref setup
% ------------------------------------------------------------------
\hypersetup{
    colorlinks=true,
    linkcolor=terra,
    urlcolor=teal,
    pdfauthor={Qhawarina},
    pdftitle={Reporte Trimestral Qhawarina <<QUARTER>>},
    pdfsubject={Análisis Económico del Perú},
}

% ------------------------------------------------------------------
% Header / footer (standard pages)
% ------------------------------------------------------------------
\pagestyle{fancy}
\fancyhf{}
\renewcommand{\headrulewidth}{0.4pt}
\renewcommand{\headrule}{\hbox to\headwidth{\color{border}\leaders\hrule height \headrulewidth\hfill}}
\fancyhead[L]{\small\color{ink3}QHAWARINA — Reporte Trimestral <<QUARTER>>}
\fancyhead[R]{\small\color{ink3}Análisis Económico del Perú}
\fancyfoot[L]{\footnotesize\color{ink3}Qhawarina}
\fancyfoot[C]{\footnotesize\color{ink3}\thepage}
\fancyfoot[R]{\footnotesize\color{ink3}Generado el <<DATE>>}

% ------------------------------------------------------------------
% Section styling
% ------------------------------------------------------------------
\titleformat{\section}[block]
  {\normalfont\large\bfseries\color{terra}}
  {\thesection.}{0.5em}{}[\vspace{-0.3em}\textcolor{border}{\hrule height 0.5pt}]

\titleformat{\subsection}[block]
  {\normalfont\normalsize\bfseries\color{ink}}
  {\thesubsection.}{0.4em}{}

\titlespacing*{\section}{0pt}{1.4em}{0.7em}
\titlespacing*{\subsection}{0pt}{0.9em}{0.4em}

% ------------------------------------------------------------------
% tcolorbox styles
% ------------------------------------------------------------------
\tcbuselibrary{skins, breakable}

\newtcolorbox{notabox}[1]{
  enhanced,
  colback=bg,
  colframe=terra,
  boxrule=0pt,
  leftrule=3pt,
  arc=0pt,
  left=8pt, right=8pt, top=6pt, bottom=6pt,
  fonttitle=\bfseries\color{terra},
  title=#1,
}

\newtcolorbox{kpibox}[1]{
  enhanced,
  colback=white,
  colframe=border,
  boxrule=0.5pt,
  arc=2pt,
  left=8pt, right=8pt, top=6pt, bottom=6pt,
  fonttitle=\small\bfseries\color{ink3},
  title=#1,
  attach boxed title to top left={yshift=-2mm, xshift=6pt},
  boxed title style={
    colback=bg,
    colframe=border,
    boxrule=0.5pt,
    arc=1pt,
  },
}

\newtcolorbox{sectionbox}[1]{
  enhanced,
  colback=white,
  colframe=border,
  boxrule=0.5pt,
  arc=3pt,
  left=10pt, right=10pt, top=8pt, bottom=8pt,
  fonttitle=\bfseries\color{terra},
  title=#1,
}

% ------------------------------------------------------------------
% Custom commands
% ------------------------------------------------------------------
\newcommand{\kpivalue}[1]{%
  {\fontsize{16}{18}\selectfont\bfseries\color{ink}#1}%
}
\newcommand{\kpilabel}[1]{%
  {\small\color{ink3}#1}%
}
\newcommand{\quarter}{<<QUARTER>>}

% ============================================================
\begin{document}

% ============================================================
% COVER PAGE
% ============================================================
\thispagestyle{empty}
\begin{center}

\vspace*{2cm}

\begin{tcolorbox}[
  enhanced,
  colback=terra,
  colframe=terra,
  boxrule=0pt,
  arc=4pt,
  left=20pt, right=20pt, top=20pt, bottom=20pt,
  width=0.85\textwidth,
  halign=center,
]
  {\fontsize{32}{36}\selectfont\bfseries\color{white}QHAWARINA}\\[6pt]
  {\large\color{white!80!terra}Análisis Económico del Perú}
\end{tcolorbox}

\vspace{1.5cm}

{\fontsize{24}{28}\selectfont\bfseries\color{ink}Reporte Trimestral}\\[0.5em]
{\fontsize{20}{24}\selectfont\color{terra}\textbf{<<QUARTER>>}}

\vspace{0.8em}

{\large\color{ink3}<<QUARTER_START>> — <<QUARTER_END>>}

\vspace{2cm}

\begin{tcolorbox}[
  enhanced,
  colback=white,
  colframe=border,
  boxrule=0.5pt,
  arc=3pt,
  left=15pt, right=15pt, top=10pt, bottom=10pt,
  width=0.75\textwidth,
]
  \centering
  \small\color{ink3}
  Este reporte ha sido generado automáticamente por el pipeline\\
  Qhawarina a partir de datos públicos del BCRP, INEI, MIDAGRI\\
  y fuentes de supermercados peruanos.\\[6pt]
  \textbf{Fecha de generación:} <<DATE>>
\end{tcolorbox}

\vspace{1.5cm}

% KPI snapshot on cover
\noindent
\begin{minipage}[t]{0.22\textwidth}
  \begin{kpibox}{BPP}
    \centering
    \kpivalue{<<BPP_YOY>>}\\[2pt]
    \kpilabel{Var. acum.}
  \end{kpibox}
\end{minipage}%
\hfill
\begin{minipage}[t]{0.22\textwidth}
  \begin{kpibox}{Riesgo Pol.}
    \centering
    \kpivalue{<<RISK_SCORE>>}\\[2pt]
    \kpilabel{/100}
  \end{kpibox}
\end{minipage}%
\hfill
\begin{minipage}[t]{0.22\textwidth}
  \begin{kpibox}{PBI}
    \centering
    \kpivalue{<<GDP_NOWCAST>>}\\[2pt]
    \kpilabel{nowcast i.a.}
  \end{kpibox}
\end{minipage}%
\hfill
\begin{minipage}[t]{0.22\textwidth}
  \begin{kpibox}{Pobreza}
    \centering
    \kpivalue{<<POVERTY_RATE>>}\\[2pt]
    \kpilabel{nacional}
  \end{kpibox}
\end{minipage}

\end{center}

% ============================================================
\newpage
% ============================================================
% TABLE OF CONTENTS
% ============================================================
\thispagestyle{plain}
\tableofcontents
\vspace{1em}

\begin{tcolorbox}[
  enhanced,
  colback=bg,
  colframe=border,
  boxrule=0.5pt,
  arc=2pt,
  left=6pt, right=6pt, top=4pt, bottom=4pt,
]
  {\footnotesize\color{ink3}
    Reporte generado automáticamente. Todos los nowcasts son estimaciones estadísticas.
    Los datos de supermercados cubren Plaza Vea, Metro y Wong (Lima Metropolitana).
  }
\end{tcolorbox}

% ============================================================
\newpage
% ============================================================
% SECTION 1 — RESUMEN EJECUTIVO
% ============================================================
\pagestyle{fancy}

\section{Resumen Ejecutivo}

\begin{notabox}{PERSPECTIVA TRIMESTRAL}
  \small\color{ink}
  <<QUARTERLY_PERSPECTIVA>>
\end{notabox}

\vspace{0.8em}

\subsection{Cuadro Resumen de Indicadores}

\begin{center}
\begin{tabular}{lrrl}
  \toprule
  \textbf{Indicador} & \textbf{Inicio trimestre} & \textbf{Fin trimestre} & \textbf{Variación} \\
  \midrule
<<ALL_KPI_TABLE>>  \bottomrule
\end{tabular}
\end{center}

\vspace{0.8em}

\subsection{Eventos Relevantes del Trimestre}

\begin{description}[leftmargin=3.5cm, labelindent=0pt, style=sameline]
<<EVENTS_LIST>>
\end{description}

% ============================================================
\newpage
% ============================================================
% SECTION 2 — PRECIOS Y CANASTA BPP
% ============================================================

\section{Precios y Canasta BPP}

\subsection{Evolución del Índice BPP}

\begin{center}
  \includegraphics[width=0.90\textwidth]{<<CHART_BPP_QUARTER>>}
\end{center}

\vspace{0.6em}

\subsection{Variación por Categoría}

\begin{center}
\begin{tabular}{lrrll}
  \toprule
  \textbf{Categoría} & \textbf{Var. trimestral} & \textbf{Var. interanual} & \textbf{Trim. anterior} & \textbf{Aceleración} \\
  \midrule
<<CATEGORY_QUARTER_TABLE>>  \bottomrule
\end{tabular}
\end{center}

\vspace{0.6em}

\begin{tcolorbox}[
  enhanced,
  colback=bg,
  colframe=border,
  boxrule=0.5pt,
  arc=2pt,
  left=6pt, right=6pt, top=4pt, bottom=4pt,
]
  {\footnotesize\color{ink3}
    \textbf{Metodología BPP:} Índice Jevons bilateral, cadena diaria.
    Fuentes: Plaza Vea, Metro, Wong. Filtro de ratio extremo: $0.5 < r < 2.0$.
    Base 100 = inicio del período de muestra. Peso IPC: estructura INEI Lima dic. 2021.
  }
\end{tcolorbox}

% ============================================================
\newpage
% ============================================================
% SECTION 3 — PBI Y MACROECONOMÍA
% ============================================================

\section{PBI y Macroeconomía}

\subsection{Track Record del Nowcast}

\begin{center}
\begin{tabular}{lrrl}
  \toprule
  \textbf{Trimestre} & \textbf{Nowcast} & \textbf{Dato oficial} & \textbf{Error} \\
  \midrule
<<GDP_TRACK_RECORD_TABLE>>  \bottomrule
\end{tabular}
\end{center}

\vspace{0.6em}

\begin{tcolorbox}[
  enhanced,
  colback=bg,
  colframe=border,
  boxrule=0.5pt,
  arc=2pt,
  left=6pt, right=6pt, top=4pt, bottom=4pt,
]
  {\footnotesize\color{ink3}
    Nowcast generado con DFM ($k=3$ factores) y bridge equation Ridge ($\alpha=1$).
    Ventana rodante 7 años. Series: 33--37 del panel nacional BCRP/INEI.
    RMSE histórico: $\sim$5.4pp (backtest 2004--2024, excluyendo COVID).
    Pre-COVID RMSE: $\sim$1.4pp.
  }
\end{tcolorbox}

\vspace{0.8em}

\subsection{Mercado Cambiario}

\begin{center}
  \includegraphics[width=0.85\textwidth]{<<CHART_FX_QUARTER>>}
\end{center}

\vspace{0.4em}

\noindent
\begin{minipage}[t]{0.49\textwidth}
  \begin{kpibox}{Tipo de Cambio}
    \kpivalue{S/ <<FX_RATE>>}\\[2pt]
    \kpilabel{PEN/USD — cierre del trimestre}\\[4pt]
    \kpilabel{Cambio reciente: <<FX_DAILY_CHANGE>>}
  \end{kpibox}
\end{minipage}%
\hfill
\begin{minipage}[t]{0.49\textwidth}
  \begin{kpibox}{Política Monetaria}
    \kpivalue{<<BCRP_RATE>>}\\[2pt]
    \kpilabel{Tasa de referencia BCRP}\\[4pt]
    \kpilabel{<<FX_QUARTER_SUMMARY>>}
  \end{kpibox}
\end{minipage}

% ============================================================
\newpage
% ============================================================
% SECTION 4 — RIESGO POLÍTICO
% ============================================================

\section{Riesgo Político}

\subsection{Evolución del Índice de Inestabilidad}

\begin{center}
  \includegraphics[width=0.90\textwidth]{<<CHART_RISK_QUARTER>>}
\end{center}

\vspace{0.6em}

\noindent
\begin{minipage}[t]{0.49\textwidth}
  \begin{kpibox}{Nivel Actual}
    \kpivalue{<<RISK_SCORE>>/100}\\[2pt]
    \kpilabel{Nivel: \textbf{<<RISK_LEVEL>>}}\\[4pt]
    \kpilabel{Artículos monitoreados: <<RISK_ARTICLES>>}
  \end{kpibox}
\end{minipage}%
\hfill
\begin{minipage}[t]{0.49\textwidth}
  \begin{tcolorbox}[
    enhanced,
    colback=white,
    colframe=border,
    boxrule=0.5pt,
    arc=2pt,
    left=8pt, right=8pt, top=6pt, bottom=6pt,
  ]
    {\small\color{ink3}
      Zonas de referencia:\\
      \textcolor{teal}{$\bullet$} Bajo: 0--30\\
      \textcolor{amber}{$\bullet$} Moderado: 30--60\\
      \textcolor{red}{$\bullet$} Alto: 60--100\\[4pt]
      Fuente: Clasificación NLP sobre 6+ fuentes RSS.
    }
  \end{tcolorbox}
\end{minipage}

% ============================================================
\newpage
% ============================================================
% SECTION 5 — POBREZA DEPARTAMENTAL
% ============================================================

\section{Pobreza Departamental}

\subsection{Nowcast Departamental 2025}

\begin{center}
  \includegraphics[width=0.85\textwidth]{<<CHART_POVERTY_DEPT>>}
\end{center}

\vspace{0.6em}

\subsection{Tabla: Departamentos por Tasa de Pobreza}

\begin{center}
\begin{tabular}{lrrr}
  \toprule
  \textbf{Departamento} & \textbf{Nowcast 2025} & \textbf{Dato 2024} & \textbf{Cambio} \\
  \midrule
<<DEPT_POVERTY_TABLE>>  \bottomrule
\end{tabular}
\end{center}

\vspace{0.4em}

\begin{tcolorbox}[
  enhanced,
  colback=bg,
  colframe=border,
  boxrule=0.5pt,
  arc=2pt,
  left=6pt, right=6pt, top=4pt, bottom=4pt,
]
  {\footnotesize\color{ink3}
    \textbf{Metodología:} Gradient Boosting Regressor con features rodantes (rolling 12 meses).
    Panel: 26 departamentos, 2004--2024. Lag de publicación ENAHO: 18 meses.
    RMSE: $\sim$2.5pp; Rel.RMSE vs AR1: 0.95. IC 95\% basado en errores históricos.
    Callao incluido en Lima para compatibilidad con ENAHO histórico.
  }
\end{tcolorbox}

% ============================================================
\newpage
% ============================================================
% SECTION 6 — NOTAS METODOLÓGICAS
% ============================================================

\section{Notas Metodológicas}

\subsection{Cambios Metodológicos este Trimestre}

\begin{notabox}{CAMBIOS METODOLÓGICOS}
  \small\color{ink}
  <<METHODOLOGY_CHANGES>>
\end{notabox}

\vspace{0.8em}

\subsection{Modelos y Fuentes de Datos}

\begin{sectionbox}{Panel de Datos Nacional}
  \small\color{ink}
  El panel nacional mensual contiene 58 series en formato largo con las siguientes variables:
  \texttt{date}, \texttt{series\_id}, \texttt{series\_name}, \texttt{category},
  \texttt{value\_raw}, \texttt{value\_sa}, \texttt{value\_log}, \texttt{value\_dlog},
  \texttt{value\_yoy}, \texttt{source}, \texttt{frequency\_original}, \texttt{publication\_lag\_days}.

  \vspace{4pt}
  \textbf{Fuentes principales:} BCRP (API pública), INEI (publicaciones mensuales),
  MIDAGRI (boletines mayoristas y avícolas), Supermercados (scraper propio),
  NASA VIIRS (Night-Time Lights departamental).
\end{sectionbox}

\vspace{0.6em}

\begin{tcolorbox}[
  enhanced,
  colback=white,
  colframe=border,
  boxrule=0.5pt,
  arc=3pt,
  left=10pt, right=10pt, top=8pt, bottom=8pt,
]
  \small\color{ink}
  \textbf{Nowcast PBI (DFM + Ridge Bridge)}\\
  Modelo de Factores Dinámicos con $k=3$ factores no observados, ventana rodante de 7 años.
  Bridge equation Ridge ($\alpha=1.0$) mapea factores a variación trimestral del PBI.
  Exclusión de COVID (2020--2021) en entrenamiento. RMSE pre-COVID: $\approx 1.4$pp.\\[8pt]

  \textbf{Nowcast Inflación (DFM)}\\
  Objetivo: variación mensual del IPC (media móvil 3 meses). 19--23 series candidatas.
  Rel.RMSE vs AR(1): $\approx 0.99$. Incluye precios mayoristas, pollos, supermercados.\\[8pt]

  \textbf{Nowcast Pobreza (GBR)}\\
  Gradient Boosting con features: PBI departamental, inflación regional, NTL, rezagos de pobreza.
  Entrenamiento: panel 2004--2024 (excl. COVID). RMSE $\approx 2.5$pp, mejor que AR(1).\\[8pt]

  \textbf{Riesgo Político (NLP)}\\
  Clasificación bayesiana sobre 81+ feeds RSS. Score normalizado con componente financiero
  (FX, spread crédito, reservas). Índice diario desde 2020.\\[8pt]

  \textbf{Índice BPP}\\
  Jevons bilateral, cadena diaria, $\sim$37,000--42,000 productos observados.
  Cobertura: arroz, pan, carnes, pescados, lácteos, frutas, verduras, aceites,
  azúcar, bebidas, limpieza, cuidado personal.
\end{tcolorbox}

\vspace{1em}

\begin{tcolorbox}[
  enhanced,
  colback=bg,
  colframe=border,
  boxrule=0.5pt,
  arc=2pt,
  left=6pt, right=6pt, top=4pt, bottom=4pt,
]
  {\footnotesize\color{ink3}
    \textbf{Qhawarina} (\textit{quechua}: "el que observa") es una plataforma de análisis
    económico en tiempo real para el Perú, desarrollada con metodologías de punta en
    nowcasting macroeconómico. Los resultados son de carácter informativo y no constituyen
    pronóstico oficial del Banco Central de Reserva del Perú (BCRP), el Ministerio de
    Economía y Finanzas (MEF) ni el Instituto Nacional de Estadística e Informática (INEI).\\[4pt]
    Código fuente y documentación: repositorio privado Nexus.
    Contacto: pipeline@qhawarina.pe
  }
\end{tcolorbox}

\end{document}
